\documentclass[bigger,pdf]{beamer}
\usepackage[utf8]{inputenc}
\usepackage[T1]{fontenc}
\usepackage{graphicx}
\usepackage{amsmath}
\usepackage{amssymb}
\usepackage{amsfonts}
\usepackage{algorithm,algorithmic}

\include{preamble}

%Informações do Slide Inicial
\title{Computer Networks}
\subtitle{Lecture 1}
\author{Lucas Dias Hiera Sampaio}
%Para o institute sempre atentar para a ordem de serviço 05/2019 de 28 de agosto de 2019
\institute{Universidade Tecnológica Federal do Paraná (UTFPR),\\ Câmpus Cornélio Procópio, Programa de Pós Graduação em Informática} 

\begin{document}

\begin{frame}
    \titlepage
\end{frame}

\section{Overview}
\begin{frame}[plain]
    \frametitle{\contentsname}
    %https://tex.stackexchange.com/questions/109748/
    \begin{columns}[t]
        \begin{column}{.5\textwidth}
            \tableofcontents[sections={1-6}]
        \end{column}
        \begin{column}{.5\textwidth}
            \tableofcontents[sections={7-12}]
        \end{column}
    \end{columns}
\end{frame}

\section{Introduction}

\begin{frame}{Introduction}

\begin{block}{Pythagoras Theorem}
Let $a$ be the hypotenuse of a right triangle, $b$ and $c$ its \textit{catheti} or legs, then:
\begin{equation}
    a^2 = b^2 \,+\, c^2
\end{equation}
\end{block}

\end{frame}

\section{Blocks}

\begin{frame}{Blocks}

\begin{block}{Block Example}
This is a simple block example.
\end{block}
\begin{alertblock}{Alert Example}
This is a simple alert block example.
\end{alertblock}
\begin{exampleblock}{Example Example}
This is a simple example block example
\end{exampleblock}

\end{frame}

\section{Equations}

\begin{frame}{Equations}
Equation (\ref{eq:2}) defines the theoretic channel capacity given by the Shannon-Hartley Theorem:
\begin{equation}\label{eq:2}
    C \, = \, B\,\log_2\left(1+\delta\right)
\end{equation}
\noindent where $B$ is the channel bandwidth and $\delta$ is the signal-to-noise ratio (SNR).
\end{frame}

\subsection{Multiple Equations}

\begin{frame}{Multiple Equations}
\begin{eqnarray}
    (a+b)^2 & = & 0 \\
    a^2 + 2ab + b^2 & = & 0 \nonumber \\
    a^2 + b^2 &=& -2ab 
\end{eqnarray}
\end{frame}

\subsection{Matrices}

\begin{frame}{Matrices}
\begin{eqnarray} M &=& \left[ {\begin{array}{cccc}
    m_{1,1} & m_{1,2} & \ldots & m_{1,N} \\
    m_{2,1} & m_{2,2} & \ldots & m_{1,N} \\
    \vdots & \vdots & \ddots & \vdots \\
    m_{N,1} & m_{N,2} & \ldots & m_{N,N}
\end{array}} \right] \\ \nonumber
\end{eqnarray}
\end{frame}

\section{Figures}

\begin{frame}{Figures}
\begin{block}{How to Include Figures}
There are basically two ways to include figures in a beamer presentation:
\begin{itemize}
    \item When there are little to no details that are small or when the figure size does not matter, one may include it inside a frame as shown in the next slide.
    \item When there are many details and the figure must be enlarged one may use a full frame to show the figure and ignore frame default content as shown next.
\end{itemize}
\end{block}
\end{frame}

\begin{frame}{Figures inside}
\begin{figure}
    \centering
    \includegraphics[width=.5\textwidth]{figures/figure1.png}
    \caption{Example of a Figure inside a frame.}
    \label{fig:1}
\end{figure}
\end{frame}

%\begin{frame}{Figures in Full Screen}
{
\setbeamertemplate{navigation symbols}{}
\begin{frame}[plain]
    \makebox[\linewidth]{\includegraphics[width=\paperwidth]{figures/figure1.png}}
\end{frame}
}
%\end{frame}

\section{Tables}

\begin{frame}{Tables}
\begin{table}[]
    \centering
    \begin{tabular}{|c|c|}
        \hline
        A & B \\
        \hline
        C & D \\
        \hline
    \end{tabular}
    \caption{Table Example.}
    \label{tab:1}
\end{table}
\end{frame}

\section{Algorithms}

\begin{frame}{Algorithms}
\begin{algorithm}[H]
\begin{algorithmic}[1]
\FOR{$i=1$ to $N$}
\FOR{$j=1$ to $JJJJ$}
\STATE $energy[i*JJJ+j] =$ 
$ interpolate(AAA[i*JJJ+j], ZZZ)$
\ENDFOR
\ENDFOR
\end{algorithmic}
\caption{pseudocode for the calculation of }
\label{alg:seq}
\end{algorithm}
\end{frame}

\section{Code Examples}

\begin{frame}[fragile]{Code Examples}

\begin{verbatim}
    int main() {
        printf("Hello World");
        return 0;
    }
\end{verbatim}

\end{frame}

\section{References}

\begin{frame}{References}
  Some references to showcase [allowframebreaks] \cite{knuth92,ConcreteMath,Simpson,Er01,greenwade93}
\end{frame}

\begin{frame}[allowframebreaks]{References}

  \bibliography{refs}
  \bibliographystyle{abbrv}

\end{frame}



\end{document}